\documentclass{article}
\usepackage{tabularx}
\usepackage{booktabs}

\begin{document}

\begin{table}[htbp]
  \centering
  \caption{DNA Storage vs. Traditional Storage}
  \label{tab:comparison}
  \begin{tabularx}{\textwidth}{llX}
    \toprule
    \textbf{Criteria} & \textbf{DNA Storage} & \textbf{Traditional Storage} \\
    \midrule
    Storage Density & High capacity; immense data storage potential. & Limited capacity; constrained by physical size. \\
    Durability & Long-term stability; potential for centuries. & Susceptible to degradation; limited lifespan. \\
    Longevity & Potentially indefinite; under optimal conditions. & Limited lifespan; requires periodic migration. \\
    Data Retrieval & Time-consuming; advanced sequencing required. & Quick access; immediate data retrieval. \\
    Environmental Impact & Minimal; compact storage reduces footprint. & Significant; requires energy-intensive facilities. \\
    Cost & High initial costs; decreasing with advancements. & Variable; influenced by technology and capacity. \\
    Security & Potential for enhanced encryption and security. & Established encryption methods; potential vulnerabilities. \\
    Archival Potential & Suitable for long-term archival; data preservation. & Limited archival; risk of data loss over time. \\
    Scalability & Scalable; potential for exponential growth. & Limited scalability; physical constraints. \\
    Application & Emerging technology; research and development phase. & Established technology; widespread use. \\
    \bottomrule
  \end{tabularx}
\end{table}

\end{document}